% Especificacao de Caso de Uso: sucGerenciarSolicitacoes
%
% Conteudo convertido do documento original do NPI. A tipografia segue o
% relatorio de estagio (unifil.cls): Helvetica em tudo, 12pt, A4, margens
% 3cm/2cm, paragrafo com recuo de 1,5cm e sumario com pontilhado.
% Do original ficaram a capa, o cabecalho em tabela de 3 linhas e o rodape
% "Confidencial", que sao a identidade do documento do NPI.

\documentclass[12pt,a4paper]{article}

\usepackage[utf8]{inputenc}
\usepackage[T1]{fontenc}
\usepackage[brazilian]{babel}
\usepackage[inner=30mm,outer=20mm,top=20mm,bottom=20mm,includehead,
            headheight=62pt,headsep=10mm,footskip=12mm]{geometry}
\usepackage{graphicx}
\usepackage{float}
\usepackage{indentfirst}
\usepackage{setspace}
\usepackage{helvet}
\usepackage{array}
\usepackage{tabularx}
\usepackage{fancyhdr}
\usepackage{titlesec}
\usepackage{tocloft}
\usepackage{enumitem}
\usepackage[font=normalsize,labelfont=bf,labelsep=endash,
            justification=centering,position=above]{caption}
\usepackage[colorlinks=true,linkcolor=black,urlcolor=black,citecolor=black]{hyperref}

\renewcommand{\familydefault}{\sfdefault}
\onehalfspacing
\setlength{\parindent}{1.5cm}

% Folga extra na justificacao: os nomes dos casos de uso nao hifenizam e
% sem ela uma linha que contenha dois deles estoura a margem.
\setlength{\emergencystretch}{4em}

% Os nomes dos casos de uso nao devem ser hifenizados.
\hyphenation{sucGerenciarUsuarios sucGerenciarAmbientes sucGerenciarPredios
             sucGerenciarSalas sucGerenciarTiposDeSalas sucGerenciarRecursos
             sucGerenciarAgendamentos sucGerarRelatorios sucExportarRelatorios
             sucGerenciarSolicitacoes saProfessor saCoordenador}

% ------------------------------------------------------------- metadados

\newcommand{\ucnome}{sucGerenciarSolicitacoes}
\newcommand{\ucprojeto}{Reserva de salas}
\newcommand{\ucversao}{1.0}
\newcommand{\ucdata}{30/10/2026}
\newcommand{\ucautor}{José Ronaldo dos Santos Jr.}
\newcommand{\ucorgao}{Núcleo de Prática de Informática - NPI}
\newcommand{\uctitulo}{Especificação de Caso de Uso: \ucnome}

% --------------------------------------------------- cabecalho e rodape

\pagestyle{fancy}
\fancyhf{}
\renewcommand{\headrulewidth}{0pt}
\renewcommand{\footrulewidth}{0pt}

\newcolumntype{H}[1]{>{\raggedright\arraybackslash}p{#1}}

\fancyhead[L]{%
  \setstretch{1.0}\small
  \begin{tabularx}{\textwidth}{|X|H{4.4cm}|}
    \hline
    \textbf{\ucprojeto} & Versão:\qquad \ucversao \\
    \hline
    \uctitulo & Data:\ \ \ucdata \\
    \hline
    \uctitulo{} - \ucautor & \\
    \hline
  \end{tabularx}%
}

\fancyfoot[L]{\small Confidencial}
\fancyfoot[C]{\small\copyright\ucorgao, 2026}
\fancyfoot[R]{\small Página \thepage}

% ------------------------------------------------------ titulos de secao

% "1." nas secoes e "3.1" nas subsecoes, como no documento do NPI.
\renewcommand{\thesection}{\arabic{section}.}
\renewcommand{\thesubsection}{\arabic{section}.\arabic{subsection}}

\titleformat{\section}{\normalfont\large\bfseries}{\thesection}{0.8em}{}
\titleformat{\subsection}{\normalfont\normalsize\bfseries}{\thesubsection}{0.8em}{}
\titlespacing*{\section}{0pt}{1.4\baselineskip}{0.5\baselineskip}
\titlespacing*{\subsection}{0pt}{1.0\baselineskip}{0.4\baselineskip}

% ------------------------------------------------------------- o sumario

% Mesmo formato do relatorio: titulo centralizado em caixa alta, numeros a
% esquerda, pontilhado ate a pagina, negrito no primeiro nivel.
\addto\captionsbrazilian{\renewcommand{\contentsname}{SUMÁRIO}}

\renewcommand{\cfttoctitlefont}{\hfill\normalfont\Large}
\renewcommand{\cftaftertoctitle}{\hfill}
\setlength{\cftbeforetoctitleskip}{0pt}
\setlength{\cftaftertoctitleskip}{1.5\baselineskip}

\renewcommand{\cftsecfont}{\bfseries}
\renewcommand{\cftsecpagefont}{\bfseries}
\renewcommand{\cftsecleader}{\bfseries\cftdotfill{\cftdotsep}}
\renewcommand{\cftsubsecfont}{\normalfont}
\renewcommand{\cftsubsecpagefont}{\normalfont}

% Os dois niveis compartilham a mesma coluna de numeros e de titulos.
\setlength{\cftsecindent}{0pt}
\setlength{\cftsecnumwidth}{1.8cm}
\setlength{\cftsubsecindent}{0pt}
\setlength{\cftsubsecnumwidth}{1.8cm}
\setlength{\cftbeforesecskip}{0.7\baselineskip}
\setlength{\cftbeforesubsecskip}{0pt}

% ---------------------------------------------------------- os fluxos

\newlist{fluxo}{enumerate}{1}
\setlist[fluxo]{label=\arabic*., leftmargin=1.5cm, labelsep=0.5em,
                topsep=0.4\baselineskip, itemsep=0.2\baselineskip, parsep=0pt}

\begin{document}

% ------------------------------------------------------------------ capa

\thispagestyle{empty}
% geometry[includehead] reserva o espaco do cabecalho tambem nesta pagina,
% que nao tem cabecalho; o recuo negativo devolve o topo.
\vspace*{-3.3cm}
{\bfseries
  \hrule height 0.8pt
  \vspace{10pt}
  \raggedleft \Large \ucorgao \par
  \vspace{10pt}
  \hrule height 0.8pt
  \vspace{36pt}
  \raggedleft \Large \ucprojeto \par
  \vspace{12pt}
  \raggedleft \Large \uctitulo \par
  \vspace{36pt}
  \raggedleft \large Versão \ucversao \par
}
\newpage

% -------------------------------------------- historico da revisao e sumario

\begin{center}
  {\Large\bfseries Histórico da Revisão}
\end{center}

\vspace{\baselineskip}

{\setstretch{1.0}\small
\noindent
\begin{tabularx}{\textwidth}{|p{2.6cm}|p{1.6cm}|X|p{3.6cm}|}
  \hline
  \centering\textbf{Data} & \centering\textbf{Versão} & \centering\textbf{Descrição}
    & \centering\arraybackslash\textbf{Autor} \\
  \hline
  \ucdata & \ucversao & Primeira versão & \ucautor \\
  \hline
  & & & \\
  \hline
  & & & \\
  \hline
  & & & \\
  \hline
\end{tabularx}
\par}

\vspace{\baselineskip}

% O sumario do relatorio e entrelinhado simples, com respiro so entre os
% grupos de primeiro nivel.
{\setstretch{1.0}\tableofcontents\par}
% \tableofcontents forca \thispagestyle{plain}; devolvemos cabecalho e rodape.
\thispagestyle{fancy}
\newpage

% ------------------------------------------------------------- o documento

\begin{center}
  {\Large\bfseries \uctitulo}
\end{center}

\vspace{0.5\baselineskip}

O objetivo deste documento é apresentar todas as especificações dos requisitos
funcionais do caso de uso \ucnome.

\section{Breve Descrição}

A função deste caso de uso é permitir que o usuário do sistema, coordenador ou
professor, seja capaz de visualizar, submeter, aprovar, rejeitar e reverter
solicitações de alteração de agendamentos, com suas devidas permissões no
sistema.

\section{Fluxo Básico de Eventos}

\begin{fluxo}
  \item O usuário acessa o menu principal.
  \item O usuário acessa a seção de solicitações.
  \item O usuário escolhe a opção de submeter uma solicitação.
  \item O sistema exibe um formulário com as alterações pretendidas e a
        justificativa.
  \item O usuário faz a inserção dos dados.
  \item O usuário submete a solicitação.
  \item O sistema valida os dados recebidos, registra a solicitação como
        pendente e informa os choques de sala previstos.
\end{fluxo}

\section{Fluxos Alternativos}

Nesta seção, serão abordados fluxos alternativos relacionados ao caso de uso
\ucnome. Um fluxo alternativo representa um comportamento alternativo,
geralmente derivado de exceções que ocorrem no fluxo principal.

\subsection{Visualizar Solicitações}

\begin{fluxo}
  \item O usuário acessa o menu principal.
  \item O usuário acessa a seção de solicitações.
  \item O usuário seleciona uma solicitação.
  \item O usuário visualiza as informações desejadas.
\end{fluxo}

\subsection{Aprovar Solicitações}

\begin{fluxo}
  \item O usuário acessa o menu principal.
  \item O usuário acessa a seção de solicitações.
  \item O usuário seleciona a solicitação pendente que deseja decidir.
  \item O usuário seleciona a opção de aprovar e registra um comentário.
  \item O sistema efetiva as alterações nos agendamentos e finaliza o processo.
\end{fluxo}

\subsection{Rejeitar Solicitações}

\begin{fluxo}
  \item O usuário acessa o menu principal.
  \item O usuário acessa a seção de solicitações.
  \item O usuário seleciona a solicitação pendente que deseja decidir.
  \item O usuário seleciona a opção de rejeitar e registra um comentário.
  \item O sistema mantém os agendamentos inalterados e finaliza o processo.
\end{fluxo}

\subsection{Reverter Solicitações}

\begin{fluxo}
  \item O usuário acessa o menu principal.
  \item O usuário acessa a seção de solicitações.
  \item O usuário seleciona uma solicitação já decidida.
  \item O usuário seleciona a opção de reverter e informa a justificativa.
  \item O sistema abre uma nova solicitação pendente com as operações
        contrárias.
\end{fluxo}

\subsection{Tratamento de Dados}

O sistema valida os dados inseridos pelo usuário. O sistema apresenta uma
mensagem de erro caso os dados inseridos pelo usuário não estejam adequados
para o sistema.

\section{Subfluxos}

Não existem subfluxos associados a este caso de uso.

\section{Cenários Chave}

Os cenários-chave deste caso de uso são: submissão, visualização, aprovação,
rejeição e reversão de solicitações registradas no sistema.

\section{Condições Prévias}

Não se aplica.

\section{Condições Posteriores}

Nesta seção serão listados os possíveis estados do sistema imediatamente após
a conclusão do caso de uso em questão.

\subsection{Submeter Solicitação}

Uma solicitação será registrada no sistema com o status pendente e salva no
banco de dados.

\subsection{Aprovar Solicitação}

Uma solicitação pendente passará para o status aprovado e as alterações serão
aplicadas aos agendamentos no banco de dados.

\subsection{Rejeitar Solicitação}

Uma solicitação pendente passará para o status rejeitado e os agendamentos
permanecerão inalterados no banco de dados.

\subsection{Reverter Solicitação}

Uma nova solicitação será registrada no banco de dados com o status pendente,
apontando para a solicitação original.

\section{Pontos de Extensão}

O caso de uso \ucnome{} não possui pontos de extensão.

\section{Requisitos Especiais}

O caso de uso \ucnome{} não possui requisitos especiais.

\section{Informações Adicionais}

Na Figura~\ref{fig:caso-de-uso} são apresentados os atores e funcionalidades
do caso de uso.

\begin{figure}[H]
  \caption{Caso de uso}
  \label{fig:caso-de-uso}
  \centering
  \includegraphics[width=0.6\textwidth]{../../../out/src/plantuml/casos-de-uso/sucGerenciarSolicitacoes.png}

  {\footnotesize \textbf{Fonte:} O autor.}
\end{figure}

\end{document}
